\documentclass[11pt,a4paper]{article}
\usepackage[margin=1in]{geometry}
\usepackage{amsmath, mathtools, amssymb, amsthm}
\usepackage{graphicx}
\usepackage{booktabs}
\usepackage{xcolor}
\usepackage[colorlinks=true, linkcolor=blue, citecolor=blue, urlcolor=blue]{hyperref}

\theoremstyle{definition}
\newtheorem{definition}{Definition}
\newtheorem{remark}{Remark}

\newcommand{\todo}[1]{\vspace{0.3em}\noindent\fbox{\parbox{0.95\textwidth}{\textit{\textbf{Draft note:} #1}}}\vspace{0.3em}}
\newcommand{\IR}{\mathrm{IR}}
\newcommand{\EUR}{EUR\,}

\title{\textbf{From Dead Cross-Sectional Momentum to Belief-State Deep\\
Time-Series Momentum:\\
An Honest, Cost-Aware Study on a Diversified Asset Basket}}
\author{Stefanos Drakos\\
\small AGEL AI I.K.E., Rhodes, Greece\\
\small ORCID: 0000-0001-7417-2444}
\date{\today}

\begin{document}
\maketitle

\begin{abstract}
We study momentum trading through the lens of honest, cost-aware, out-of-sample evaluation,
and report a deliberately two-sided result. First, a \emph{negative}: on US large-cap equities
(S\&P~100/500), cross-sectional momentum---ranking names against one another---is dead net of
realistic costs over 2006--2024, and a Bayesian online changepoint overlay does not revive it.
Second, a \emph{positive}: the same trading logic applied as \emph{time-series} momentum to a
diversified basket of eighteen liquid exchange-traded funds (equities, bonds, gold, commodities,
the US dollar, real estate, sectors) is alive. We construct a per-asset belief state with a
Kalman local-linear-trend filter and Bayesian online changepoint detection, feed it together with
multi-horizon trend features to a small LSTM ``Deep Momentum Network'' trained on a portfolio
Sharpe-ratio objective net of turnover costs, and select hyperparameters by a \emph{nested}
walk-forward that never lets the validation choice touch the out-of-sample test. The learned
model attains an out-of-sample (2016--2024) information ratio of $0.58$ (Newey--West
$t\approx 1.98$, Deflated Sharpe Ratio $0.67$ over five trials), beating a fixed multi-horizon
rules baseline ($0.47$) at lower drawdown. We isolate the source of the edge: it is
\emph{asset-class diversity}, not name count---five genuinely uncorrelated instruments
(average pairwise correlation $0.03$) more than double the risk-adjusted return of twenty-two
sector-diversified equities (correlation $0.41$). Finally we show the practical use: a
return-stacked overlay (full equities plus a half-sized trend sleeve) beats buy-and-hold S\&P~500
out of sample on both absolute return and Sharpe. All experiments are reproducible from free Yahoo
data, net of a transparent cost model, with pre-registered honesty: a correctly-inferred null is
reported as such.
\end{abstract}

\section{Introduction}
\label{sec:intro}

Momentum---the tendency of recent winners to keep winning---is among the most studied phenomena in
empirical finance, and one of the few with out-of-sample persistence across decades and asset
classes \cite{moskowitz2012}. The deep-learning literature has enriched it: Deep Momentum Networks
inject learned trading rules into a volatility-scaling framework and optimize a Sharpe-ratio loss
directly \cite{lim2019}; Slow Momentum with Fast Reversion augments this with changepoint detection
to survive trend reversals \cite{wood2021}; and the Momentum Transformer replaces recurrence with
attention \cite{wood2021mt}. A recurring weakness of the broader reinforcement- and deep-learning
trading literature is evaluation: raw price windows as state, in-sample tuning, and performance
reported without rigorous walk-forward or cost accounting \cite{hambly2023}.

This paper takes the honest-evaluation stance seriously and lets it dictate the result rather than
the other way around. We make three contributions, organized as a single empirical narrative.

\textbf{(i) A motivating negative.} We first show that the obvious adaptation---cross-sectional
momentum on US large-cap equities, dollar-neutral long-winners/short-losers---is \emph{dead} net
of costs in 2006--2024, and that a Bayesian online changepoint overlay does not save it
(Section~\ref{sec:negative}). This is consistent with the decay of equity momentum in liquid
large-caps and is reported plainly, not buried.

\textbf{(ii) A belief-state Deep Momentum Network that works on the right universe.} We construct a
per-asset belief state from a Kalman local-linear-trend filter (filtered velocity, a trend
$t$-statistic, and the standardized innovation) and Bayesian online changepoint detection
\cite{kalman1960,adams2007}, feed it with multi-horizon volatility-normalized returns to a small
shared LSTM, and train on a portfolio Sharpe objective net of turnover costs in the direct-
reinforcement spirit of \cite{moody2001}. Crucially, hyperparameters are chosen by a \emph{nested}
walk-forward---validation-only selection with early stopping---so the out-of-sample number is
honest by construction (Section~\ref{sec:method}). On a diversified eighteen-ETF basket the learned
model beats a strong fixed-rule baseline out of sample (Section~\ref{sec:results}).

\textbf{(iii) Where the edge comes from, and how to use it.} We isolate the mechanism---asset-class
diversity, quantified by average pairwise correlation---and show it dominates name count
(Section~\ref{sec:diversity}); and we demonstrate a practical deployment in which a return-stacked
trend overlay beats buy-and-hold equity out of sample (Section~\ref{sec:deploy}).

We are deliberately conservative about magnitudes. Better state estimation extracts a weak signal
more cleanly; it does not manufacture signal. Under the Fundamental Law of Active Management
\cite{grinold1999}, our levers are signal quality and breadth of \emph{independent} bets, not raw
instrument count. The realistic target is a diversified information ratio of order $0.4$--$0.7$,
and every number below is reported net of a transparent cost model.

\section{Related Work}
\label{sec:related}

\paragraph{Momentum and deep momentum.} Time-series momentum \cite{moskowitz2012} trades each asset
on its own past return, volatility-scaled; it is the economic baseline our model must beat. Deep
Momentum Networks \cite{lim2019} learn position sizing end-to-end under a Sharpe loss; the Slow
Momentum / Fast Reversion architecture \cite{wood2021} adds changepoint detection to handle
turning points; the Momentum Transformer \cite{wood2021mt} is the attention-based successor. Our
model is a deliberately small LSTM Deep Momentum Network whose novelty is not capacity but the
\emph{belief-state} inputs (a single Kalman filter supplies both the trend and the changepoint
signal) and the \emph{nested}, cost-aware evaluation.

\paragraph{Filtering and changepoints.} The Kalman filter \cite{kalman1960} and structural
time-series models give the linear-Gaussian belief; Bayesian online changepoint detection
\cite{adams2007} supplies a per-asset regime-break severity. We note that simple linear models are
competitive with elaborate sequence models at this frequency and signal-to-noise
\cite{zeng2023}, which motivates a small network plus an informative inductive bias rather than raw
capacity.

\paragraph{Honest evaluation.} Multiple-testing inflation of backtested Sharpe ratios is corrected
by the Deflated Sharpe Ratio \cite{bailey2014}; walk-forward discipline, purged cross-validation,
and the probability of backtest overfitting are developed in \cite{lopezdeprado2018}. We use
Newey--West $t$-statistics throughout and report the Deflated Sharpe Ratio against the number of
configurations tried.

\section{A Motivating Negative: Cross-Sectional Equity Momentum}
\label{sec:negative}

We first test the natural equity adaptation: each day, rank the investable universe by a
$12{-}1$ momentum score, go long the top decile and short the bottom decile, dollar-neutral and
volatility-scaled, rebalanced and held, net of a $5$~bps spread on turnover. On the S\&P~100
(2000--2024) and on a $414$-name S\&P~500 sample (2006--2024), every variant we tried lost money
out of sample: plain cross-sectional momentum, a Kalman-velocity ranking, and a
significance-gated belief variant all returned information ratios between $-0.16$ and $-1.34$ on
the S\&P~100 and between $-0.28$ and $-0.47$ on the broader sample across daily to quarterly
rebalancing. A market-level Bayesian changepoint overlay did not help (it moved the information
ratio from $-0.36$ to $-0.37$): the losses are not concentrated in detectable regime breaks but
bleed across a structurally low-momentum, highly-correlated, efficient large-cap universe. We
report this as a clean null and move to the setting where the same logic is alive.

\section{Belief-State Features and the Deep Momentum Network}
\label{sec:method}

\paragraph{Universe.} A fixed basket of eighteen liquid ETFs spanning asset classes: US large/
tech/small-cap, developed and emerging international equity; long, intermediate, investment-grade
and high-yield bonds; gold, silver, broad commodities and oil; the US dollar; real estate; and
three equity sectors. Daily adjusted prices from Yahoo Finance, 2007--2024.

\paragraph{Per-asset belief state.} For each asset we run a Kalman local-linear-trend filter on the
log price \cite{kalman1960}, reading off the filtered trend velocity $\hat v_t$, the trend
significance $\hat v_t/\sqrt{P^{vv}_t}$, and the standardized innovation $\nu_t/\sqrt{S_t}$; and a
Bayesian online changepoint detector \cite{adams2007} on the returns, summarized as a break
severity $s_t\in[0,1]$ (posterior mass on short run lengths). The feature vector per asset per day
($F=10$) is: five volatility-normalized multi-horizon returns ($1,21,63,126,252$ days), log
realized volatility, and the four belief signals $(\hat v_t,\ \hat v_t/\sqrt{P^{vv}_t},\
\nu_t/\sqrt{S_t},\ s_t)$.

\paragraph{Model and objective.} A single shared LSTM reads each asset's own feature sequence and
emits a position $\pi_{j,t}=\tanh(\cdot)\in[-1,1]$; there is no cross-asset information at decision
time. The training objective is the Sharpe ratio of the equal-capital portfolio net of turnover
cost,
\begin{equation}
\mathcal{L} \;=\; -\,\frac{\bar R}{\mathrm{std}(R)}, \qquad
R_t = \frac{1}{N}\sum_{j=1}^N \pi_{j,t}\,\rho_{j,t+1} \;-\; c\,\frac{1}{N}\sum_j |\pi_{j,t}-\pi_{j,t-1}|,
\label{eq:sharpe}
\end{equation}
with next-day returns $\rho$ and cost rate $c$, in the direct-reinforcement spirit of
\cite{moody2001}. Diversification is thus inside the loss, even though each position is decided
from a single asset's features.

\paragraph{Honest nested walk-forward.} We use two expanding out-of-sample folds (test blocks
ending 2020 and 2024). Within each fold, the last $20\%$ of the training window is a validation
split used \emph{both} for early stopping and for selecting the configuration (hidden size,
weight decay, dropout) from a small grid; the test block is touched exactly once, for measurement.
Feature standardization is fit on training data only. This makes the reported out-of-sample number
free of selection bias by construction.

\section{Results}
\label{sec:results}

Table~\ref{tab:main} reports the headline out-of-sample comparison (2016--2024, net of costs,
each return stream scaled to $10\%$ annualized volatility for comparability). The learned Deep
Momentum Network attains an information ratio of $0.58$ (Newey--West $t=1.98$; Deflated Sharpe
Ratio $0.67$ against five configurations), beating the fixed multi-horizon rules baseline ($0.47$)
at a smaller maximum drawdown ($-17\%$ vs $-21\%$). The nested procedure selected the same model
size (hidden $=16$) in both folds, indicating a stable rather than fragile choice.

\begin{table}[h]
\centering
\begin{tabular}{lrrrr}
\toprule
Strategy (out-of-sample 2016--2024, net) & IR & NW-$t$ & DSR & max DD \\
\midrule
Rules time-series momentum (fixed)            & $0.47$ & --     & --     & $-21\%$ \\
ML LSTM Deep Momentum Network (nested)        & $\mathbf{0.58}$ & $1.98$ & $0.67$ & $\mathbf{-17\%}$ \\
\bottomrule
\end{tabular}
\caption{Honest nested out-of-sample comparison on the eighteen-ETF basket. Figures
\texttt{etf\_ml\_vs\_rules.png}.}
\label{tab:main}
\end{table}

\section{Where the Edge Comes From: Diversity, Not Count}
\label{sec:diversity}

The Fundamental Law of Active Management \cite{grinold1999} ties realized information ratio to the
number of \emph{independent} bets, not the raw instrument count. We make this concrete. A
fixed time-series-momentum rule on twenty-two sector-diversified US equities returns a Sharpe ratio
of only $0.15$, with average pairwise correlation $0.41$: across eleven sectors the names still
fall together, so the effective breadth is small. Five genuinely uncorrelated diversifiers (long
and intermediate Treasuries, gold, broad commodities, the dollar) have average pairwise correlation
$0.03$ and a standalone Sharpe of $0.19$. Combining the two---twenty-seven instruments---yields a
Sharpe of $0.34$, \emph{higher than either component alone}: the diversification ``free lunch.''
Consistently, five well-chosen diversified ETFs outperform one hundred correlated equities
(information ratio $+0.54$ vs $+0.09$) under the identical rule. The lesson is that asset-class
diversity, measured by low cross-correlation, is the lever---not the number of names or sectors.

\section{Practical Deployment: Beating Buy-and-Hold by Combining}
\label{sec:deploy}

Over 2009--2024 the US equity market was an exceptional bull; in such a regime no diversified,
risk-controlled strategy beats $100\%$ equities on \emph{absolute} return, and ours does not in
isolation. The honest and useful question is how the trend sleeve \emph{combines} with equities.
Out of sample (2016--2024, net), a return-stacked overlay---hold the full equity position and add a
half-sized trend sleeve via margin ($1.5\times$ gross)---turns \EUR$10{,}000$ into \EUR$37{,}390$
versus \EUR$29{,}825$ for buy-and-hold S\&P~500, with a higher Sharpe ($0.98$ vs $0.84$) at
comparable drawdown; a leverage-free $50/50$ blend has the best risk-adjusted profile (Sharpe
$1.04$, maximum drawdown roughly halved). The mechanism is that the trend sleeve earns precisely
when equities fall (it was positive through the 2022 selloff, short bonds and equities, long the
dollar and energy), so stacking it on equities adds return where it hurts most. Figure
\texttt{etf\_beat\_buyhold.png} shows the equity curves; we flag the obvious risk that the
overlay uses leverage and that the equity--trend diversification, while historically negative in
crises, is not guaranteed.

\section{Scope and Limitations}
\label{sec:scope}

The result is conditional on the tested universe and period; survivorship and the specific ETF set
are caveats, and daily data preclude microstructure effects. The learned edge ($\IR\approx 0.58$,
$t\approx 2$) is positive but not overwhelming, as is honest for nine years of out-of-sample
diversified data. The Kalman belief is linear-Gaussian by design; genuine fat tails and stochastic
volatility are only partially captured. Costs are a flat spread plus, where shorting equities,
a financing charge; true execution slippage and borrow are not modeled. Consistent with prior
practice, where evidence is weak we say so: the cross-sectional equity result is a null, and the
time-series result is a modest, cost-surviving, but not spectacular edge.

\section{Conclusion}
\label{sec:concl}

Honest evaluation, not architecture, is the protagonist. The same momentum logic is dead on
cross-sectional large-cap equities and alive on a diversified asset basket; a belief-state Deep
Momentum Network, selected by a leak-free nested walk-forward, beats a strong fixed rule out of
sample; the edge traces to asset-class diversity rather than instrument count; and the practical
payoff is a return-stacked overlay that beats buy-and-hold equity. Every claim is reproducible from
free data, net of costs, with nulls reported as nulls.

\begin{thebibliography}{99}

\bibitem{adams2007} Adams, R.\ P., \& MacKay, D.\ J.\ C.\ (2007). Bayesian online changepoint
detection. \emph{arXiv preprint} arXiv:0710.3742.

\bibitem{bailey2014} Bailey, D.\ H., \& L\'opez de Prado, M.\ (2014). The Deflated Sharpe Ratio:
correcting for selection bias, backtest overfitting, and non-normality. \emph{The Journal of
Portfolio Management} 40(5):94--107.
\href{https://doi.org/10.3905/jpm.2014.40.5.094}{doi:10.3905/jpm.2014.40.5.094}

\bibitem{grinold1999} Grinold, R.\ C., \& Kahn, R.\ N.\ (1999). \emph{Active Portfolio Management},
2nd ed. McGraw-Hill.

\bibitem{hambly2023} Hambly, B., Xu, R., \& Yang, H.\ (2023). Recent advances in reinforcement
learning in finance. \emph{Mathematical Finance} 33(3):437--503.

\bibitem{kalman1960} Kalman, R.\ E.\ (1960). A new approach to linear filtering and prediction
problems. \emph{Journal of Basic Engineering} 82(1):35--45.

\bibitem{lim2019} Lim, B., Zohren, S., \& Roberts, S.\ (2019). Enhancing time series momentum
strategies using deep neural networks. \emph{arXiv preprint} arXiv:1904.04912. (\emph{The Journal
of Financial Data Science}, 2019.)

\bibitem{lopezdeprado2018} L\'opez de Prado, M.\ (2018). \emph{Advances in Financial Machine
Learning}. Wiley.

\bibitem{moody2001} Moody, J., \& Saffell, M.\ (2001). Learning to trade via direct reinforcement.
\emph{IEEE Transactions on Neural Networks} 12(4):875--889.
\href{https://doi.org/10.1109/72.935097}{doi:10.1109/72.935097}

\bibitem{moskowitz2012} Moskowitz, T.\ J., Ooi, Y.\ H., \& Pedersen, L.\ H.\ (2012). Time series
momentum. \emph{Journal of Financial Economics} 104(2):228--250.

\bibitem{wood2021} Wood, K., Roberts, S., \& Zohren, S.\ (2021). Slow momentum with fast reversion:
a trading strategy using deep learning and changepoint detection. \emph{arXiv preprint}
arXiv:2105.13727. (\emph{The Journal of Financial Data Science}, 2022.)

\bibitem{wood2021mt} Wood, K., Giegerich, S., Roberts, S., \& Zohren, S.\ (2021). Trading with the
momentum transformer: an intelligent and interpretable architecture. \emph{arXiv preprint}
arXiv:2112.08534.

\bibitem{zeng2023} Zeng, A., Chen, M., Zhang, L., \& Xu, Q.\ (2023). Are transformers effective for
time series forecasting? \emph{Proceedings of the AAAI Conference on Artificial Intelligence}
37(9):11121--11128. \href{https://doi.org/10.1609/aaai.v37i9.26317}{doi:10.1609/aaai.v37i9.26317}

\end{thebibliography}

\vspace{0.5em}
\noindent\textit{Reproducibility.} All code, the cost model, and the figures are in
\texttt{etoro/paper4/} and \texttt{etoro/Strategies/slow-momentum-fast-reversion/}; results
regenerate from free Yahoo data via \texttt{paper4/code/run\_etf.py}.

\end{document}
